% !Mode:: "TeX:UTF-8"
% !Mode:: "TeX:UTF-8"
\newcommand{\thepage}{} % to silence a hyperref warning about pdfpagelabels
\RequirePackage[]{silence}
\WarningsOff[hyperref]
\PassOptionsToPackage{log-declarations=false}{xparse}
%\RequirePackage[l2tabu, orthodox]{nag}
\RequirePackage{luatex85} % to make \pdf...= commands work with LuaTeX
\documentclass[aspectratio=169,t]{beamer}
\usepackage{ifthen}
\usepackage[T1]{fontenc}
\usepackage{iftex}
\ifLuaTeX
 \usepackage{fontspec}
 \setmainfont[Ligatures=TeX]{Source Sans Pro}
\else
 \usepackage[utf8]{inputenc}
\fi
\usepackage{german}
\usepackage{newunicodechar}
\newunicodechar{€}{\texteuro}
\newunicodechar{ß}{\ss}
\renewcommand{\SS}{\iffontchar\font"1E9E \symbol{"1E9E}\else SS\fi}
\usepackage{ellipsis}
\usepackage{microtype}
\usepackage{listings}
\usepackage{sourcecodepro}
\lstset{language=Python, basicstyle=\ttfamily}
\usepackage{amsfonts}
\usepackage{amssymb}
\usepackage{amsthm}
\usepackage{fancyhdr}
\usepackage{graphicx}

\newtheorem{proposition}[theorem]{Proposition}

\DeclareMathOperator*{\argmax}{argmax}
\DeclareMathOperator{\supp}{supp}
\newcommand{\E}{\ensuremath{\mathbb{E}}}
\newcommand{\V}{\ensuremath{\mathbb{V}}}

\setbeamercolor{emph}{fg=purple}
\renewcommand<>{\emph}[1]{%
  {\usebeamercolor[fg]{emph}\only#2{\itshape}#1}%
}



\hypersetup{draft=false,colorlinks=false,pdfpagemode=UseOutlines,plainpages=false,
            breaklinks=true,bookmarksnumbered=true,pdfborder=0 0 0,
            pdftitle={Data Science SoSe 2026 HHU},
            pdfauthor={Konrad Voelkel},
            pdfsubject={data science}}

%\documentclass[aspectratio=169,t]{beamer}
\usetheme{hhubeamer} % put the hhu-beamer-theme in /usr/share/texmf/tex/latex/beamer/themes/

% -------------------- Code listings config --------------------
% Define custom colors for syntax highlighting
\definecolor{codegreen}{rgb}{0,0.6,0}
\definecolor{codegray}{rgb}{0.5,0.5,0.5}
\definecolor{codepurple}{rgb}{0.58,0,0.82}
\definecolor{backcolour}{rgb}{0.96,0.96,0.94} % Very light gray/beige
% Configure the listings package
\lstset{
    language=Python,
    backgroundcolor=\color{backcolour},   
    commentstyle=\color{codegreen}\itshape,
    keywordstyle=\color{blue}\bfseries,
    numberstyle=\tiny\color{codegray},
    stringstyle=\color{codepurple},
    basicstyle=\ttfamily\footnotesize, % Use a monospaced font
    breakatwhitespace=false,         
    breaklines=true,      
             % Enable automatic line breaking
    captionpos=b,                      % Caption at the bottom
    keepspaces=true,                 
    numbers=left,                      % Line numbers on the left
    numbersep=10pt,                    % Distance of line numbers from the code
    showspaces=false,                
    showstringspaces=false,            % Don't show weird underscores in strings
    showtabs=false,                  
    tabsize=4,
    frame=single,                      % Add a subtle border around the code
    rulecolor=\color{black!20},        % Border color
  }

\event[]{}
\author[Sommersemester 2026]{Dr. Konrad Völkel, Sommersemester 2026}

%\title{Data Science, VL 1}
%\date{16.04.2026}


%%% Local Variables:
%%% mode: latex
%%% TeX-master: t
%%% End:


\title{Data Science, VL 3} % neu in 2026
\date{23.04.2025}

\begin{document}

\begin{frame}
  \titlepage
\end{frame}


\section{Einführung in Python}

\begin{frame}
  \frametitle{Einführung in Python}
  \begin{itemize}[<+->]
  \item REPL ist unser Freund
  \item Einrückung statt Klammern
  \item Alles ist ein Objekt
  \item LLMs können Python
  \end{itemize}
  \bigskip
  \href{https://media.hhu.de/channel/video/data-science-python-in-der-repl/1383cd1ccc36f42479354b1c8c0dd373/167}{Erste Schritte in der REPL (Video) unter media.hhu.de/channel/video/data-science-python-in-der-repl/1383cd1ccc36f42479354b1c8c0dd373/167}
\end{frame}

\subsection{Datentypen}

\begin{frame}
  \frametitle{Zahlen: \texttt{int} und \texttt{float}}
  \begin{itemize}[<+->]
  \item \texttt{int}: Ganze Zahlen, die beliebig groß werden können (kein Overflow).
  \item \texttt{float}: Kommazahlen (IEEE 754 binary64 double precision auf den meisten Plattformen).
  \item Typen müssen nicht explizit deklariert werden, Python bindet Variablen dynamisch.
  \item Operatoren: \texttt{+}, \texttt{-}, \texttt{*}, \texttt{/}.
  \item Besonderheiten: \texttt{//} für ganzzahlige Division und \texttt{\%} für Modulo.
  \end{itemize}
\end{frame}

\begin{frame}
  \frametitle{Text: Strings (\texttt{str})}
  \begin{itemize}[<+->]
  \item \texttt{str} repräsentiert Text als Sequenzen von Unicode-Zeichen.
  \item Strings werden mit einfachen oder doppelten Anführungszeichen erstellt.
  \item Strings sind unveränderlich (immutable); Modifikationen erzeugen neue Objekte.
  \item Zugriff auf einzelne Zeichen oder Teilbereiche erfolgt über Slicing (\texttt{x[1:4]}).
  \end{itemize}
\end{frame}

\begin{frame}
  \frametitle{Wahrheitswerte und \texttt{None}}
  \begin{itemize}[<+->]
  \item \texttt{bool}: Der Datentyp für Wahrheitswerte mit den Konstruktoren \texttt{True} und \texttt{False}.
  \item Nahezu jedes Objekt kann in Python zu einem Boolean evaluiert werden (z.B. leere Sammlungen sind \texttt{False}).
  \item \texttt{None}: Ein spezielles Objekt (null in Java).
  \item Nützlich, um Variablen zu deklarieren, die noch keinen Wert haben.
  \end{itemize}
\end{frame}

\begin{frame}
  \frametitle{Listen (\texttt{list})}
  \begin{itemize}[<+->]
  \item Eine geordnete, veränderbare und dynamische Sammlung von Objekten (ähnlich ArrayList).
  \item Werden mit eckigen Klammern \texttt{[]} oder \texttt{list()} initialisiert.
  \item Längenabfrage funktioniert einfach über die eingebaute Funktion \texttt{len(liste)}.
  \end{itemize}
\end{frame}

\begin{frame}
  \frametitle{Dictionaries (\texttt{dict})}
  \begin{itemize}[<+->]
  \item Speichern von Key-Value-Paaren, (eine Hash-Map).
  \item Die Syntax nutzt geschweifte Klammern und Doppelpunkte:\\
    \texttt{\{\dq key\dq: \dq value\dq\}}.
  \item Als Keys (Schlüssel) sind nur unveränderliche Typen wie\\ \texttt{str}, \texttt{int} oder \texttt{float} erlaubt.
  \item Wenn ein angefragter Key nicht existiert, wirft Python einen \texttt{KeyError}.
  \end{itemize}
\end{frame}

\subsection{Kontrollstrukturen}

\begin{frame}
  \frametitle{Fallunterscheidungen (\texttt{if} / \texttt{else})}
  \begin{itemize}[<+->]
  \item Python verzichtet auf Klammern zur Abgrenzung von Codeblöcken.
  \item Blöcke werden durch einen Doppelpunkt eingeleitet und eingerückt\\ (4 Leerzeichen sind 'ein Tab').
  \end{itemize}
\end{frame}

\begin{frame}
  \frametitle{Schleifen (\texttt{while} und \texttt{for})}
  \begin{itemize}[<+->]
  \item \texttt{while}-Schleifen prüfen eine Bedingung und wiederholen den Körper, bis diese nicht mehr zutrifft.
  \item \texttt{for}-Schleifen in Python sind eigentlich for-each-Schleifen: Sie iterieren direkt über Sammlungen wie Listen oder Dictionaries.
  \item Um klassisch über Zahlen zu iterieren, nutzt man das Objekt \texttt{range(a, b)}.
  \end{itemize}
\end{frame}

\subsection{Funktionen \& I/O}

\begin{frame}
  \frametitle{Funktionen definieren (\texttt{def})}
  \begin{itemize}[<+->]
  \item Funktionen werden mit dem Schlüsselwort \texttt{def} eingeleitet.
  \item Im Gegensatz zu \texttt{if} oder Schleifen erzeugt \texttt{def} einen eigenen Geltungsbereich (Scope) für Variablen.
  \item Funktionen können optionale Argumente (mit Default-Werten) und Keyword-Argumente annehmen.
  \end{itemize}
\end{frame}

\begin{frame}[separator]
  \frametitle{Vielen Dank für Ihre Aufmerksamkeit!}
\end{frame}

\end{document}

%%% Local Variables:
%%% mode: latex
%%% TeX-master: t
%%% End:
